\chapter{速查表与最后冲刺}

\section{最常用 Java 片段}

\begin{javacode}
// 邻接表
List<List<Integer>> graph = new ArrayList<>();
for (int i = 0; i < n; i++) graph.add(new ArrayList<>());

// 访问标记
boolean[] visited = new boolean[n];

// BFS
Deque<Integer> queue = new ArrayDeque<>();
queue.offer(start);
while (!queue.isEmpty()) {
    int current = queue.poll();
}

// 最小堆 / 最大堆
PriorityQueue<Integer> minHeap = new PriorityQueue<>();
PriorityQueue<Integer> maxHeap =
    new PriorityQueue<>(Comparator.reverseOrder());

// 哈希表默认容器
Map<Integer, List<Integer>> groups = new HashMap<>();
groups.computeIfAbsent(key, ignored -> new ArrayList<>()).add(value);

// 排序
Arrays.sort(nums);
Arrays.sort(points, (a, b) -> Integer.compare(a[0], b[0]));
\end{javacode}

\section{C++ 到 Java 的高频替换}

\begin{center}
\begin{tabularx}{\linewidth}{@{}l X@{}}
\toprule
当你想写 C++ & Java 中先想 \\
\midrule
\code{nullptr} & \code{null}，并先决定在哪里检查 \\
\code{new Node(...)} + \code{delete} & \code{new Node(...)}；只维护可达关系 \\
\code{vector<int> dp} & \code{int[] dp} \\
\code{unordered_map[key]} & \code{getOrDefault} 或 \code{computeIfAbsent} \\
\code{priority_queue<int>} & \code{PriorityQueue<Integer>}（最小堆）\\
\code{greater<int>} & \code{Comparator.reverseOrder()}（用于最大堆）\\
\code{sort(v.begin(),v.end())} & \code{Arrays.sort(array)} 或 \code{list.sort(...)} \\
\code{INT_MIN} & \code{Integer.MIN_VALUE} \\
\code{long long} & \code{long} \\
\bottomrule
\end{tabularx}
\end{center}

\section{你应该保留的 C++ 优点}

\begin{itemize}
  \item 写代码前先说清楚复杂度和不变量；
  \item 图、树、堆题先定义状态，再动手；
  \item 对边界、空指针、哨兵和溢出保持敏感；
  \item 手写 LRU、Treap 的经历保留了数据结构直觉，不要因为 Java 有库就丢掉。
\end{itemize}

\begin{transfer}[毕业标准]
当你看到 ``图 + 依赖关系'' 时，第一反应是拓扑排序；看到 ``流数据中位数'' 时，第一反应是双堆；看到 Java 时，第一反应也应当是相同的算法，只是顺手拿起 \code{ArrayDeque}、\code{HashMap}、\code{PriorityQueue} 和 \code{StringBuilder}。
\end{transfer}

\begin{drill}[最终模拟]
限时 45 分钟，用 Java 完成以下三题：课程表、LRU Cache、数据流中位数。完成后逐项执行第六章的 90 秒检查。若你卡住的是 API，就重做本章速查；若你卡住的是状态设计，回到对应 C++ 解法先读不变量，不要急着看 Java 答案。
\end{drill}
